\chapter[\emj{🫧}~Ordenamientos Básicos (Bubble, Insertion, Selection)]{\emj{🫧}~Ordenamientos Básicos (Bubble, Insertion, Selection)}

\begin{dsaquees}

Tres formas manuales de ordenar cosas en la vida real:

1. Bubble Sort 🫧 (Burbujas): Imagina burbujas subiendo en un refresco.
   Comparas dos elementos juntos; si el de la izquierda es mayor, lo
   cambias. Los números más grandes "burbujean" hacia el final en 
   cada pasada.

2. Insertion Sort 🃏 (Cartas): Como cuando ordenas cartas en tu mano.
   Tomas una carta nueva y la vas comparando con las que ya tienes
   ordenadas hasta que encuentras su lugar exacto y la "insertas".

3. Selection Sort 🔍 (Selección): Buscas el objeto más pequeño de 
   todo el grupo y lo pones al principio. Luego buscas el más pequeño
   de los que quedan y lo pones en segundo lugar, y así...

\end{dsaquees}

\begin{dsaotraforma}

Son algoritmos de ordenamiento $O(n^2)$. Aunque son lentos para grandes
cantidades de datos, son fundamentales para entender la lógica de 
sorting.

- Bubble Sort: Se basa en intercambios adyacentes.
- Insertion Sort: Divide el array en una parte ordenada y otra no. 
  Es muy eficiente para arrays que ya están casi ordenados.
- Selection Sort: Se basa en encontrar el mínimo en cada iteración.

\end{dsaotraforma}

\begin{dsadiagrama}
\begin{verbatim}
Bubble Sort (Paso 1):
[64, 34, 25, 12] -> 64 > 34? SÍ -> [34, 64, 25, 12]
[34, 64, 25, 12] -> 64 > 25? SÍ -> [34, 25, 64, 12]
[34, 25, 64, 12] -> 64 > 12? SÍ -> [34, 25, 12, 64]* (64 ya está al final)

Insertion Sort (Paso 2 - ordenando el 25):
Mano: [34, 64] | Pendiente: [25, 12]
Tomo el 25. ¿25 < 64? SÍ. ¿25 < 34? SÍ.
Mano: [25, 34, 64] | Pendiente: [12]

Selection Sort (Paso 1):
[64, 34, 25, 12] -> Mínimo es 12.
Intercambio 64 y 12 -> [12, 34, 25, 64]
\end{verbatim}
\end{dsadiagrama}

\begin{lstlisting}
Bubble: 
- Compara pares (i, i+1).
- Swapea si están en orden incorrecto.
- Repite n-1 veces.

Insertion:
- Empieza en el segundo elemento.
- Compáralo hacia atrás y muévelo hasta que sea mayor que el anterior.

Selection:
- Encuentra el índice del valor mínimo en el rango [i...n].
- Intercambia `arr[i]` con `arr[min_idx]`.

PSEUDOCÓDIGO
──────────────────────────────────────────────────────────────

// Bubble Sort
PARA i desde 0 hasta n-1:
    PARA j desde 0 hasta n-i-2:
        SI arr[j] > arr[j+1]:
            intercambiar(arr[j], arr[j+1])

// Insertion Sort
PARA i desde 1 hasta n-1:
    clave = arr[i]
    j = i - 1
    MIENTRAS j >= 0 Y arr[j] > clave:
        arr[j+1] = arr[j]
        j--
    arr[j+1] = clave

CÓDIGO C++
──────────────────────────────────────────────────────────────

#include <iostream>
#include <vector>
using namespace std;

void bubbleSort(vector<int>& arr) {
    int n = arr.size();
    for (int i = 0; i < n-1; i++) {
        for (int j = 0; j < n-i-1; j++) {
            if (arr[j] > arr[j+1]) swap(arr[j], arr[j+1]);
        }
    }
}

void insertionSort(vector<int>& arr) {
    int n = arr.size();
    for (int i = 1; i < n; i++) {
        int key = arr[i];
        int j = i - 1;
        while (j >= 0 && arr[j] > key) {
            arr[j + 1] = arr[j];
            j--;
        }
        arr[j + 1] = key;
    }
}

int main() {
    vector<int> data = {64, 34, 25, 12, 22};
    insertionSort(data);
    for(int x : data) cout << x << " "; // 12 22 25 34 64
    return 0;
}

COMPLEJIDAD (Big-O)
──────────────────────────────────────────────────────────────

  Algoritmo   │ Best      │ Average   │ Worst     │ Stable?
  ────────────┼───────────┼───────────┼───────────┼─────────
  Bubble      │ O(n)      │ O(n^2)     │ O(n^2)     │ SÍ
  Insertion   │ O(n)      │ O(n^2)     │ O(n^2)     │ SÍ
  Selection   │ O(n^2)     │ O(n^2)     │ O(n^2)     │ NO

* Todos usan O(1) de espacio extra (In-place).
\end{lstlisting}

\begin{lstlisting}
✅ Insertion Sort es el mejor para arrays "casi ordenados" o para 
   pocos elementos (n < 20).

💡 Estabilidad (Stability): Significa que elementos con el mismo valor
   mantienen su orden relativo original. Bubble e Insertion son 
   estables; Selection NO.

⚠️ Casi nunca usarás estos en producción (C++ usa `std::sort` que es 
   O(n log n)), pero te los piden para evaluar tus bases.

🔑 Optimización Bubble: Puedes agregar un flag `swapped`. Si en una
   pasada completa no hubo cambios, ¡el array ya está ordenado!
\end{lstlisting}

\begin{dsaresumen}

• Los 3 son $O(n^2)$ en el peor caso.
• Bubble: El más sencillo e ineficiente.
• Insertion: El mejor para casos pequeños o casi listos.
• Selection: Hace el mínimo de swaps posibles (útil si escribir en 
  memoria es muy caro).
• In-place: No necesitan memoria extra.
════════════════════════════════════════════════════════════════

\end{dsaresumen}


